\subsection{Укороченные уравнения ГВГ}\label{subsec:shorted_shg}

Как указано в \cref{subsec:propagation}, нелинейный отклик среды описывается разложением поляризации $\mathbf{P}$ по степеням поля $\mathbf{E}$ (см.~\cref{eq:polarization}). Восприимчивость второго порядка $\chi^{(2)}$ даёт вклад в поляризацию на частоте $2\omega$:

\begin{equation}\label{eq:p_nl_chi2}
	\mathbf{P}^{(2)}(2\omega) = \varepsilon_0 \chi^{(2)} : \mathbf{E}(\omega)\mathbf{E}(\omega).
\end{equation}


\begin{equation}\label{eq:plane_wave}
	\mathbf{E}_\omega = \tfrac12 \mathbf{A}_1(z) e^{i(k_1 z - \omega t)} + \text{к.с.}, \qquad
	\mathbf{E}_{2\omega} = \tfrac12 \mathbf{A}_2(z) e^{i(k_2 z - 2\omega t)} + \text{к.с.}
\end{equation}

Для вывода укороченных уравнений ГВГ~\cite{dmitriev2004prikladnaya} представим поля основной волны и второй гармоники в виде плоских волн~\cref{eq:plane_wave} и подставим их в волновое уравнение~\cref{eq:wave} с нелинейной поляризацией~\cref{eq:p_nl_chi2}.  При выводе пользуемся приближением медленно меняющихся амплитуд, то есть пренебрегаем вторыми производными полей по координатам. Это справедливо при условии $|\partial^2 A_{1,2}/\partial z^2| \ll k_{1,2}|\partial A_{1,2}/\partial z|$. В случае ОВ необходимо заменить показатели преломления на эффективные показатели преломления соответствующих мод, а также учесть интегралы перекрытия поперечных распределений полей (см.~\cref{subsec:modes}). Итоговая система принимает вид:

\begin{subequations}\label{eq:shg_coupled}
	\begin{align}
		\frac{dA_1}{dz} &= \im \frac{2\pi\omega}{c n_\omega^{\text{eff}}} \chi^{(2)} \Gamma_1 A_1^* A_2 e^{-i\Delta k z}, \\
		\frac{dA_2}{dz} &= \im \frac{4\pi\omega}{c n_{2\omega}^{\text{eff}}} \chi^{(2)} \Gamma_2 A_1^2 e^{i\Delta k z},
	\end{align}
\end{subequations}

где $\Delta k = k_2 - 2k_1$ — волновая расстройка, а $\Gamma_{1,2}$~---интегралы перекрытия поперечных распределений полей взаимодействующих мод. В случае ГВГ в ОВ, как будет сказано далее, эти коэффициенты определяются перекрытием мод фундаментального излучения, второй гармоники и пространственного распределения наведённого статического поля $E_0$, создающего эффективную восприимчивость $\chi^{(2)} \sim \chi^{(3)} E_0$.
